\section{Experimental setup}
- TBD \subsection{Data}
\paragraph{GB1.}
The GB1 dataset is obtained from the SaProtHub repository (SaProtHub/Dataset-GB1-fitness)
and the FLIP benchmark~\cite{dallagoFLIPBenchmarkTasks2021}, using the standard two-versus-rest split
in which any variant containing at least one of the four wild-type amino acids is held out
for testing.

\paragraph{avGFP.}
The avGFP dataset is obtained from the MaveDB repository using the Sarkisyan et al.\
accession (MaveDB:00000055-a-1) and is preprocessed to retain only
variants with valid fluorescence measurements and between one and nine amino acid substitutions
relative to the wild type.
\subsection{Evaluation Protocol}

\label{sec:methods:eval}

All models are evaluated in an \emph{offline} Bayesian optimisation setting:
a fixed training set of size $n$ is sampled uniformly at random from the full dataset,
and the surrogate is used to rank the remaining (unseen) variants.
Each configuration is repeated across five independent random seeds.

Three metrics are reported: Spearman rank correlation $\rho$ between predicted and
measured fitness on the held-out set; $R^2$ on the same set; and Precision@K for
$K \in \{10, 100, 1000\}$.
Precision@K is defined as
\begin{equation}
  \text{P@}K = \frac{|\{k : \text{rank}_\text{pred}(k) \leq K\} \cap \{k : \text{rank}_\text{true}(k) \leq K\}|}{K},
\end{equation}
where $\text{rank}_\text{pred}(k)$ is the rank of variant $k$ in the predicted ordering and
$\text{rank}_\text{true}(k)$ is its rank by measured fitness.

Learning curves are generated by varying training set size
$n \in \{50, 100, 200, 500, 1000, 2000, 5000, 10000\}$ with five seeds per configuration.
\subsection{Hyperparameters and implementation details}